% ---- pg-259 （书页 247） ----
分别服从自由度为 $v_1=|(g-1)-p|+1,\ v_2=gb(n-1)-p+1$ 及 $v_1=|(b-1)-p|+1,\ v_2=gb(n-1)-p+1$ 的 $F$ 分布(见文献[1]). 在本例中,
\[
F_1=\left(\frac{1-0.3819}{0.3819}\right)\frac{(16-3+1)/2}{(|1-3|+1)/2}=7.55
\]
\[
F_2=\left(\frac{1-0.5230}{0.5230}\right)\frac{(16-3+1)/2}{(|1-3|+1)/2}=4.26
\]
且
\[
v_1=|1-3|+1=3 \qquad v_2=(16-3+1)=14
\]
由于 $F_{3,14}(0.05)=3.34$, 就有 $F_1=7.55>F_{3,14}(0.05)=3.34$. 因此我们在 5\% 水平下拒绝 $H_{01}: \boldsymbol{\tau}_1=\boldsymbol{\tau}_2=\boldsymbol{\tau}_3=\mathbf{0}$(无因子 1 效应). 类似地, $F_2=4.26>F_{3,14}(0.05)=3.34$, 因而在 5\% 显著性水平下我们拒绝 $H_{02}: \boldsymbol{\beta}_1=\boldsymbol{\beta}_2=\boldsymbol{\beta}_3=\mathbf{0}$(无因子 2 效应). 因此, 添加剂数量及挤压速度的改变对三个响应指标都有影响.

在练习 6.15 中, 研究了因子 1 和因子 2 对响应指标的影响的本质, 在该练习中, 还考虑了 $\boldsymbol{\tau}_1$ 与 $\boldsymbol{\beta}_1$ 的分量对比的联合置信区间.\hfill$\blacksquare$

\section{轮廓分析}  % 原书 6.8
\phantomsection\label{sec:6.8}

当我们将由 $p$ 种处理(如测验、问卷调查等)组成的一个处理组施加于两组或多组对象时, 可采用轮廓分析方法. 设所有响应均在相似的单元中被表示出, 再假定不同组的响应彼此独立. 我们自然要问: 总体的均值向量是否都相等? 在轮廓分析中, 诸均值向量的相等性问题可分成几种可能性.

考虑总体均值向量 $\boldsymbol{\mu}_1'=[\mu_{11},\mu_{12},\mu_{13},\mu_{14}]$, 它代表对第一组四种处理的平均响应. 这些用直线连接的均值, 其图形示于\hyperref[fig:6.4]{图 6.4}. 这种折线图就称为总体 1 的轮廓, 对每个总体(组)都可构造其轮廓. 我们主要讨论两个总体(组)的情形.

设 $\boldsymbol{\mu}_1'=[\mu_{11},\mu_{12},\cdots,\mu_{1p}]$ 和 $\boldsymbol{\mu}_2'=[\mu_{21},\mu_{22},\cdots,\mu_{2p}]$ 分别为总体 1 和总体 2 对 $p$ 种处理的平均响应. 假设 $H_0: \boldsymbol{\mu}_1=\boldsymbol{\mu}_2$ 意味着对这两个总体来说, $p$ 种处理的效应相等. 按照轮廓分析的思想方法, 可将相等性问题分成如下几步:

1. 两轮廓是否平行? 等价地说, 假设 $H_{01}: \mu_{1i}-\mu_{1,i-1}=\mu_{2i}-\mu_{2,i-1},\ i=2,3,\cdots,p$, 是否可接受?

2. 假如两轮廓是平行的, 那么它们是否重合?\footnote{练习 6.12 中考虑了这样的问题: 假如两轮廓平行, 它们是否线性? 线性轮廓的零假设可写作 $H_0: (\mu_{1i}+\mu_{2i})-(\mu_{1,i-1}+\mu_{2,i-1})=(\mu_{1,i-1}+\mu_{2,i-1})-(\mu_{1,i-2}+\mu_{2,i-2}),\ i=3,\cdots,p$. 这样的假设虽然在特定情况下人们也许会感兴趣, 但通常我们最感兴趣的是两平行轮廓是否重合, 无论它们的本质怎样.} 等价地说, 假设 $H_{02}: \mu_{1i}=\mu_{2i},\ i=1,2,\cdots,p$, 是否可以接受?

3. 假如两轮廓重合, 它们是水平的吗? 即所有均值是否都等于同一常数? 等价地说, 假设 $H_{03}: \mu_{11}=\mu_{12}=\cdots=\mu_{1p}=\mu_{21}=\mu_{22}=\cdots=\mu_{2p}$ 是否可以接受?

在第一步中, 零假设可写为
\[
H_{01}: \mathbf{C}\boldsymbol{\mu}_1=\mathbf{C}\boldsymbol{\mu}_2
\]

\begin{mvfigure}
\begin{tikzpicture}[scale=0.85]
% 坐标轴
\draw[->] (0,0)--(4.9,0) node[right,font=\small]{变量};
\draw[->] (0,0)--(0,4.5) node[above,rotate=90,font=\small]{平均响应};
% x 轴刻度
\foreach \x in {1,2,3,4} \draw (\x,0.06)--(\x,-0.06) node[below,font=\small]{\x};
% 均值水平虚线与标签
\foreach \y/\l in {1.0/\mu_{11},1.9/\mu_{12},2.8/\mu_{13}}{
  \draw[dashed] (0,\y)--(4.6,\y);
  \draw (-0.06,\y)--(0,\y);
  \node[left,font=\small] at (0,\y) {$\l$};
}
% 轮廓折线
\draw[line width=1pt] (1,1.15)--(2,3.45)--(3,2.4)--(4,4.05);
\foreach \p in {(1,1.15),(2,3.45),(3,2.4),(4,4.05)} \fill \p circle (2.4pt);
\end{tikzpicture}
\phantomsection\label{fig:6.4}\mvcap{图 6.4\quad 总体轮廓($p=4$)}
\end{mvfigure}
